\section{So sánh với các nghiên cứu liên quan}
\label{sec:so_sanh_baseline}

\subsection{Giới thiệu hệ thống mốc đối sánh của Mansoureh Lord}

Nghiên cứu của Mansoureh Lord (2021) \cite{ref_1_lord} là luận văn thạc sĩ tiêu biểu tại Sabancı University, tập trung phát hiện lỗi mất cân bằng lồng giặt trên máy giặt cửa trên sử dụng cảm biến gia tốc ADXL345 và mô hình tự mã hóa đơn giản một lớp ẩn. Hệ thống đạt độ chính xác khoảng 94,3\% ở pha vắt nhưng chưa tích hợp kênh âm thanh và chưa triển khai cơ chế chuyển đổi trạng thái tự động. Trong phần này, công trình của Mansoureh Lord được sử dụng làm hệ thống mốc đối sánh trực tiếp.

Nghiên cứu này được lựa chọn làm mốc đối sánh vì:
\begin{itemize}
    \item Cùng bài toán phát hiện bất thường trên máy giặt
    \item Sử dụng phần cứng cảm biến tương tự (ADXL345)
    \item Áp dụng kiến trúc mạng tự mã hóa cho bài toán phát hiện bất thường
    \item Công bố kết quả định lượng rõ ràng
\end{itemize}

Để xác nhận giá trị đóng góp, hệ thống tự mã hóa ba trạng thái được so sánh với nghiên cứu của Mansoureh Lord (2021) \cite{ref_1_lord} — công trình tiêu biểu ứng dụng TinyML phát hiện lỗi mất cân bằng lồng giặt.

\subsection{Cải tiến về kiến trúc mạng và thực thi phần cứng}

Sự khác biệt cốt lõi nằm ở chiến lược tối ưu hóa kiến trúc để cân bằng độ chính xác và tài nguyên:
\begin{itemize}
    \item \textbf{Cấu trúc mạng và đặc trưng đầu vào:} Hệ thống mốc đối sánh sử dụng dữ liệu rung động thô đưa trực tiếp vào mạng tự mã hóa tối giản (1 lớp ẩn) \cite{ref_1_lord}. Cách tiếp cận này nhẹ nhưng dễ bị nhiễu do không qua lọc đặc trưng. Đề tài áp dụng trích xuất đặc trưng xử lý tín hiệu số trước khi nạp vào mạng tự mã hóa Dense đối xứng. Việc tinh gọn đầu vào giúp mô hình hội tụ tốt hơn, đạt độ trễ suy luận 6,4 ms trên phần cứng nhúng — hoàn toàn đáp ứng yêu cầu thời gian thực với chu kỳ lấy mẫu 1 Hz.
    
    \item \textbf{Kỹ thuật lượng tử hóa:} Hệ thống mốc đối sánh dừng lại ở mô hình Float32 hoặc lượng tử hóa sau huấn luyện. Đề tài áp dụng lượng tử hóa nhận thức huấn luyện, nhúng sai số lượng tử hóa vào quá trình lan truyền ngược để mô hình duy trì nhất quán khi triển khai trên INT8.
\end{itemize}

\subsection{Khả năng thích nghi đa trạng thái và độ tin cậy}

Hệ thống vượt mốc đối sánh về phạm vi ứng dụng:
\begin{itemize}
    \item \textbf{Cơ chế ba trạng thái chuyên biệt:} Hệ thống mốc đối sánh tập trung vào pha vắt của máy giặt cửa trên \cite{ref_1_lord}. Hệ thống đề xuất nhận diện tự động và chuyển mô hình phù hợp với từng pha vận hành (GENTLE, STRONG, SPIN), giảm báo động giả khi thiết bị chuyển pha.
    
    \item \textbf{Kết hợp đa phương thức:} Bổ sung kênh âm thanh kỹ thuật số cho phép nhận diện các ma sát cơ khí tần số cao — điều cảm biến gia tốc đơn lẻ không thể nhận diện. Thực nghiệm cho thấy kết hợp này nâng F1-Score ở pha SPIN lên \textbf{0,9967}, vượt mốc đối sánh.
\end{itemize}

\subsection{Đối chiếu với các hướng nghiên cứu gần đây}

Ngoài mốc đối sánh TinyML trực tiếp của Lord, đồ án cũng được đặt cạnh các nghiên cứu học sâu không giám sát trên thiết bị gia dụng. Castangia và cộng sự (2022) \cite{ref_castangia_2022} sử dụng VAE trên chữ ký công suất của máy rửa bát, máy giặt và máy sấy, đạt F1-Score lớn hơn 0,90; tuy nhiên dữ liệu đầu vào của công trình này thuộc miền công suất điện, chưa xử lý trực tiếp rung động hoặc âm thanh cơ khí và chưa hướng tới triển khai TinyML trên vi điều khiển. Shul và cộng sự (2023) \cite{ref_shul_2023} khai thác phổ tiếng ồn và thông tin vận hành nội bộ của máy giặt bằng mạng DNN tự giám sát, cải thiện ROC-AUC tối đa khoảng 1,4\%; cách tiếp cận này cho thấy giá trị của tín hiệu âm thanh nhưng vẫn cần thông tin như tốc độ quay hoặc khối lượng đồ giặt từ bên trong máy.

Vì vậy, Bảng~\ref{tab:comparison_baseline} chỉ giữ phần đối chiếu trực tiếp với nghiên cứu của Lord (2021) \cite{ref_1_lord}, do đây là mốc đối sánh gần nhất với đồ án về bài toán, cảm biến và hướng triển khai TinyML. Cách trình bày này giúp tránh lặp bảng, đồng thời làm rõ đóng góp chính của hệ thống đề xuất.

\begin{table}[H]
    \centering
    \caption[Đối chiếu hiệu năng với mốc đối sánh]{Đối chiếu các thông số kỹ thuật và hiệu năng với mốc đối sánh \cite{ref_1_lord}}
    \label{tab:comparison_baseline}
    \renewcommand{\arraystretch}{1.25}
    \begin{tabularx}{\textwidth}{|>{\raggedright\arraybackslash}p{0.23\textwidth}|>{\raggedright\arraybackslash}X|>{\raggedright\arraybackslash}X|}
    \hline
    \rowcolor[HTML]{EFEFEF}
    \textbf{Chỉ số so sánh} & \textbf{Mốc đối sánh (M. Lord)} & \textbf{Hệ thống đề xuất} \\ \hline
    Đối tượng nghiên cứu & Máy giặt cửa trên & \textbf{Máy giặt cửa trước} \\ \hline
    Đặc trưng đầu vào & Rung động thô (3 trục) & \textbf{Âm thanh + rung động} \\ \hline
    Kiến trúc mô hình & Tự mã hóa Dense đơn giản & \textbf{Tự mã hóa đối xứng ba trạng thái} \\ \hline
    Lượng tử hóa & Không chuyên sâu & \textbf{INT8 toàn phần} \\ \hline
    \textbf{Độ trễ suy luận} & $> 10$ ms (ước tính) & \textbf{$6{,}4$ ms} \\ \hline
    \textbf{F1-Score (Pha SPIN)} & $\approx 0{,}943$ & \textbf{$0{,}9967$} \\ \hline
    Hệ thống cảnh báo & Đèn LED tại thiết bị & \textbf{Ứng dụng di động + đám mây thời gian thực} \\ \hline
    \end{tabularx}
\end{table}
